% main.tex
% Modular Document Example - Demonstrates \include for large projects
% KU Leuven - Faculteit Industriële Ingenieurswetenschappen
%
% This template shows how to split a large document into multiple files.
% Use this structure for thesis, large summaries, or multi-chapter documents.

\documentclass[a4paper,11pt]{report}
\usepackage{../../school-macros}

% ============================================================================
%                           DOCUMENT METADATA
% ============================================================================
\title{Warmte en Stromingen Samenvatting}
\author{Jouw Naam}
\date{Academiejaar 2025-2026}

% Optional: Define course info for headers
\newcommand{\coursename}{Warmte en Stromingen}
\newcommand{\coursecode}{}

% ============================================================================
%                              DOCUMENT
% ============================================================================
\begin{document}

% --- Title Page (KU Leuven style) ---
\kuleuventitle
    {Warmte en Stromingen Samenvatting}
    {\coursecode\ --- \coursename}
    {Jouw Naam}
    {Academiejaar 2025-2026}

\begin{abstract}
Deze samenvattingen zijn beschikbaar op GitHub waar je kunt bijdragen aan verbeteringen: \url{https://github.com/Eggmansmile/Samenvattingen-Ku-leuven-Industriele-ingenieurs}

Je bent helemaal vrij en ik raad je zelf aan om info toe te voegen aan de documenten als je de uitleg niet goed vindt.
Je leert je vak en je raakt dan ook bekend met LaTeX wat handig is voor later.
\end{abstract}

% TODO: #48 Warmte en stromingen verplaatsen naar nieuw document

% --- Front Matter ---
\tableofcontents
\newpage

% ============================================================================
%                           CHAPTER INCLUDES
% ============================================================================
% Use \include{} for chapters - each starts on a new page
% Use \input{} for smaller sections that shouldn't force page breaks
%
% Benefits of \include:
% - LaTeX can skip unchanged chapters (faster compilation)
% - Each chapter can be compiled separately for testing
% - Better organization for large projects
%
% Tip: Use \includeonly{chapter1,chapter2} in preamble to compile 
%      only specific chapters while keeping all cross-references
\printformularium
\printsymbols

\input{chapters/1_inleiding}
\chapter{Fluida}
\label{ch:fluida}
\chapter{Hydrostatica}
\label{ch:hydrostatica}

\chapter{Fluidynamics}
\label{ch:fluidynamics}
\chapter{interne stromen (viscociteit)}
\label{ch:interne_stromen_viscociteit}

\begin{enumerate}\setlength\itemsep{0pt}
	\item Viscociteit is de interne weerstand van een vloeistof tegen stroming.
	\item Er zijn twee soorten viscociteit:
\end{enumerate}


% ============================================================================
%                              APPENDICES
% ============================================================================
\appendix

% Symbol list and formularium are auto-generated

\end{document}
% --- END OF DOCUMENT ---